\documentclass[11pt]{article}
\usepackage{amsmath,amssymb,amsthm}
\usepackage[margin=1.1in]{geometry}
\usepackage{booktabs}

\newtheorem{theorem}{Theorem}
\newtheorem{definition}{Definition}
\newtheorem{remark}{Remark}

\title{Differentiation Rules:\\
The Complete Toolkit Behind Every Gradient in This Repository}
\author{Yuval Lavie}
\date{\today}

\begin{document}
\maketitle

\noindent
Every derivation in the companion notes --- activation functions, loss
functions, every MLE --- is an application of a small, closed set of
rules. This page states that set once, precisely, so the other pages can
cite rather than re-derive. Nothing here is specific to machine
learning; everything downstream is.

\section{The scalar rules}

Throughout, $f, g$ are differentiable at the point in question and
$c \in \mathbb{R}$ is a constant.

\paragraph{Linearity.} Differentiation is a linear operator:
\begin{equation}
  \frac{d}{dx}\bigl[c\,f(x) + g(x)\bigr] = c\,f'(x) + g'(x).
  \label{eq:linearity}
\end{equation}
This is why we can differentiate a sum of per-example losses term by
term --- the gradient of the empirical risk is the mean of per-example
gradients.

\paragraph{Product rule.}
\begin{equation}
  \frac{d}{dx}\bigl[f(x)\,g(x)\bigr] = f'(x)\,g(x) + f(x)\,g'(x).
  \label{eq:product}
\end{equation}

\paragraph{Quotient rule.} For $g(x) \ne 0$,
\begin{equation}
  \frac{d}{dx}\left[\frac{f(x)}{g(x)}\right]
  = \frac{f'(x)\,g(x) - f(x)\,g'(x)}{g(x)^2}.
  \label{eq:quotient}
\end{equation}
Not an independent axiom: write $f/g = f \cdot g^{-1}$ and apply the
product rule \eqref{eq:product} together with the chain rule
\eqref{eq:chain} on $g^{-1}$. The sigmoid and softmax derivations both
run on this rule.

\paragraph{Chain rule.} If $y = f(u)$ and $u = g(x)$, then
\begin{equation}
  \frac{dy}{dx} = \frac{dy}{du}\cdot\frac{du}{dx}
  = f'\bigl(g(x)\bigr)\,g'(x).
  \label{eq:chain}
\end{equation}
The single most important rule in this repository:
\emph{backpropagation is the chain rule}, applied recursively through a
composition of layers and organized to reuse intermediate products.

\paragraph{Power rule.} For any real $n$ (on the domain where $x^n$ is
defined and differentiable),
\begin{equation}
  \frac{d}{dx}\,x^n = n\,x^{n-1}.
  \label{eq:power}
\end{equation}

\paragraph{Inverse function rule.} If $f$ is invertible near $x$ with
$f'(x) \ne 0$, then
\begin{equation}
  \bigl(f^{-1}\bigr)'(y) = \frac{1}{f'\bigl(f^{-1}(y)\bigr)}.
  \label{eq:inverse}
\end{equation}
This is how the derivative of $\log$ falls out of the derivative of
$\exp$, and how the logit's derivative falls out of the sigmoid's.

\section{The transcendental derivatives}

The primitives that appear in every likelihood and every activation:
\begin{equation}
  \frac{d}{dx}\,e^{x} = e^{x},
  \qquad
  \frac{d}{dx}\,a^{x} = a^{x}\ln a,
  \qquad
  \frac{d}{dx}\,\ln x = \frac{1}{x},
  \qquad
  \frac{d}{dx}\,\log_a x = \frac{1}{x \ln a}.
  \label{eq:transcendental}
\end{equation}
And the trigonometric/hyperbolic pairs used by tanh-family activations:
\begin{equation}
  \frac{d}{dx}\sinh x = \cosh x,
  \qquad
  \frac{d}{dx}\cosh x = \sinh x,
  \qquad
  \frac{d}{dx}\tanh x = \operatorname{sech}^2 x = 1 - \tanh^2 x.
  \label{eq:hyperbolic}
\end{equation}

\paragraph{The log-derivative trick.} Combining the chain rule with
$(\ln x)' = 1/x$:
\begin{equation}
  \frac{d}{dx}\,\ln f(x) = \frac{f'(x)}{f(x)}
  \qquad\Longleftrightarrow\qquad
  f'(x) = f(x)\,\frac{d}{dx}\,\ln f(x).
  \label{eq:logderiv}
\end{equation}
This identity is why we maximize \emph{log}-likelihoods: products of
densities become sums, and the derivative of each term is a ratio. It
reappears as the score function in statistics and as the REINFORCE
estimator in reinforcement learning.

\section{Multivariable: gradient, Jacobian, and the chain rule that
backpropagation uses}

\begin{definition}[Gradient and Jacobian]
For $f : \mathbb{R}^n \to \mathbb{R}$, the gradient is the column vector
of partials, $(\nabla f)_i = \partial f / \partial x_i$. For
$\mathbf{f} : \mathbb{R}^n \to \mathbb{R}^m$, the Jacobian is the
$m \times n$ matrix
\begin{equation}
  J_{ij} = \frac{\partial f_i}{\partial x_j}.
  \label{eq:jacobian}
\end{equation}
\end{definition}

\begin{theorem}[Multivariable chain rule]
If $\mathbf{h} = \mathbf{f} \circ \mathbf{g}$ with
$\mathbf{g} : \mathbb{R}^n \to \mathbb{R}^m$ and
$\mathbf{f} : \mathbb{R}^m \to \mathbb{R}^k$, then Jacobians
\emph{multiply}:
\begin{equation}
  J_{\mathbf{h}}(x) = J_{\mathbf{f}}\bigl(\mathbf{g}(x)\bigr)\,
  J_{\mathbf{g}}(x).
  \label{eq:multichain}
\end{equation}
\end{theorem}

For a scalar loss $L$ at the end of the composition (the only case
training cares about), take the transpose of \eqref{eq:multichain} to
get the gradient recursion
\begin{equation}
  \nabla_x L = J_{\mathbf{g}}(x)^{\top}\,\nabla_{\mathbf{g}(x)} L ,
  \label{eq:backprop}
\end{equation}
i.e.\ each layer maps the \emph{upstream} gradient through its
transposed Jacobian. Backpropagation is exactly
\eqref{eq:backprop} applied layer by layer, right to left --- the
right-to-left order is what makes it cheap: every intermediate product
is a vector, never a full Jacobian.

\section{Matrix calculus identities}

Convention: gradients of scalars are column vectors shaped like the
variable (denominator layout). The identities every linear layer needs,
for $a, x \in \mathbb{R}^n$, $A \in \mathbb{R}^{n\times n}$,
$W \in \mathbb{R}^{m\times n}$:
\begin{align}
  \nabla_x \,(a^{\top}x) &= a,
  \label{eq:gradlinear}\\[2pt]
  \nabla_x \,(x^{\top}Ax) &= (A + A^{\top})\,x
  \;\;\bigl(= 2Ax \text{ for symmetric } A\bigr),
  \label{eq:gradquad}\\[2pt]
  \nabla_x \,\lVert x \rVert_2^2 &= 2x,
  \label{eq:gradnorm}\\[2pt]
  \nabla_x\,\lVert y - Wx \rVert_2^2 &= -2\,W^{\top}(y - Wx),
  \label{eq:gradlsq}\\[2pt]
  \nabla_W\,\lVert y - Wx \rVert_2^2 &= -2\,(y - Wx)\,x^{\top}.
  \label{eq:gradlsqW}
\end{align}
Equation \eqref{eq:gradlsq} is the normal-equations engine of the
linear-regression note; \eqref{eq:gradlsqW} is the weight gradient of
every dense layer under squared loss. Two derivatives of matrix
functionals that MLE for the Gaussian needs:
\begin{equation}
  \nabla_A \operatorname{tr}(BA) = B^{\top},
  \qquad
  \nabla_A \ln\det A = A^{-\top}
  \quad (A \text{ invertible}).
  \label{eq:matfunc}
\end{equation}

\section{Two rules for the boundary cases}

\paragraph{L'H\^opital.} If $f(x) \to 0$ and $g(x) \to 0$ (or both
$\to \pm\infty$) as $x \to a$, and the right-hand limit exists, then
\begin{equation}
  \lim_{x\to a} \frac{f(x)}{g(x)} = \lim_{x\to a} \frac{f'(x)}{g'(x)}.
  \label{eq:lhopital}
\end{equation}
Used to resolve $0 \cdot \log 0 = 0$ in entropy and cross-entropy, and
to show softplus $\to$ ReLU as temperature $\to 0$.

\paragraph{Subgradients.} Where a convex $f$ has a kink (ReLU at $0$,
absolute loss at $0$, hinge at margin $1$), the derivative is replaced
by the \emph{subdifferential}: the set of slopes $g$ with
$f(y) \ge f(x) + g\,(y - x)$ for all $y$. For example,
\begin{equation}
  \partial\,|x|\,\Big|_{x=0} = [-1, 1],
  \qquad
  \partial \max(0, x)\Big|_{x=0} = [0, 1].
  \label{eq:subgrad}
\end{equation}
Optimality generalizes from $f'(x) = 0$ to $0 \in \partial f(x)$;
frameworks just pick one element (PyTorch uses $0$ for ReLU at the
kink), which is harmless because the kink is hit with probability zero.

\section{Summary table}

\begin{center}
\begin{tabular}{lll}
\toprule
Rule & Statement & Where it earns its keep\\
\midrule
Linearity & $(cf+g)' = cf' + g'$ & sums of per-example losses\\
Product & $(fg)' = f'g + fg'$ & likelihood products\\
Quotient & $(f/g)' = (f'g - fg')/g^2$ & sigmoid, softmax\\
Chain & $(f\circ g)' = f'(g)\,g'$ & backpropagation\\
Power & $(x^n)' = nx^{n-1}$ & polynomial losses\\
Inverse & $(f^{-1})' = 1/f'(f^{-1})$ & logit from sigmoid\\
Log-derivative & $(\ln f)' = f'/f$ & every MLE, REINFORCE\\
Jacobian chain & $J_{f\circ g} = J_f\,J_g$ & layer composition\\
Subgradient & $0 \in \partial f$ & ReLU, hinge, absolute loss\\
\bottomrule
\end{tabular}
\end{center}

\medskip
\noindent
The companion pages apply these: \texttt{activation-functions.tex}
differentiates every activation in the repository, and
\texttt{loss-functions.tex} differentiates every loss --- each step
there cites a rule from here.

\end{document}
